% Generated draft: requires mathematical and source review.
\documentclass[11pt]{article}
\usepackage[margin=25mm]{geometry}
\usepackage{fontspec}
\setmainfont{Latin Modern Roman}
\usepackage{amsmath,amssymb,mathtools,unicode-math}
\setmathfont{Latin Modern Math}
\usepackage{enumitem,needspace,xurl,hyperref}
\setlength{\parindent}{0pt}
\setlength{\parskip}{5pt}
\setlength{\emergencystretch}{3em}
\newcommand{\N}{\mathbb{N}}
\newcommand{\Z}{\mathbb{Z}}
\newcommand{\Q}{\mathbb{Q}}
\newcommand{\R}{\mathbb{R}}
\newcommand{\C}{\mathbb{C}}
\newcommand{\F}{\mathbb{F}}
\newcommand{\A}{\mathbb{A}}
\newcommand{\PP}{\mathbb{P}}
\newcommand{\E}{\mathbb{E}}
\newcommand{\eps}{\varepsilon}
\newcommand{\dd}{\,\mathrm{d}}
\begin{document}
\section*{UnsolvedMath ID: 30000143}
\textbf{Source problem code:} catalogue entry
\subsection{Definitions and mathematical statement}
\textbf{Vanishing question, as supplied.} For every smooth scheme \(X\) over a field \(k\), must\[H^p(X,\Z(q))=0\qquad\text{for every integer }p<0\] hold, for the intended motivic weights \(q\)? The coefficient group is \(\Z\), and the degree range is strictly \(p<0\). Neither rational coefficients nor a vanishing assertion in degree \(p=0\) is substituted. The source does not explicitly state the range of \(q\), the field hypotheses, or finiteness and separatedness conditions on \(X\).\par Here \(H^p(X,\Z(q))\) denotes motivic cohomology in cohomological degree \(p\) and weight \(q\). In a usual effective-motivic convention, \(q\ge0\) and these groups are the hypercohomology of the integral motivic complex \(\Z(q)\) on the appropriate site of a smooth finite-type scheme; the particular theory and site must be fixed when extending the assertion to broader scheme classes or weights. This explanatory convention is not a source-specified replacement of the missing category. In particular, smoothness alone does not state a global finite-type or quasi-compactness restriction, and no such restriction is inserted silently into the literal question.\par The supplied title and literature triage call this Beilinson--Soul\'e vanishing. That name does not by itself fix coefficient, weight, or degree-boundary conventions: any comparison with a rational-coefficient or degree-zero formulation requires explicit hypotheses and must not change the displayed integral negative-degree question. The triage checked 2026-08-21 labels the general smooth-scheme conjecture open and reports important known special classes without identifying them or their precise theorem statements. This dated assessment requires independent verification.
\subsection{Short English statement}
The source asks for vanishing of integral motivic cohomology in negative cohomological degrees for smooth schemes over fields. The weight range and scheme category are missing. The integral coefficients and strict negative-degree range are preserved despite possible variations in the named conjecture's conventions.
\subsection{Sources}
\begin{itemize}\item Local UnsolvedMath registry, permanent ID 30000143, code catalogue entry; \emph{Beilinson--Soul\'e Vanishing for Smooth Schemes}, Oberwolfach Reports Open Problems.\item \emph{Polylogarithms (2005)}, Oberwolfach Reports, DOI 10.4171/owr/2004/48, pp. 2539--2590. \url{https://doi.org/10.4171/owr/2004/48}. Proposed year and year tag: 2004; no proposer named.\item Supplied literature label \emph{original report}: \url{https://doi.org/10.4171/owr/2004/48}.\item Supplied literature label \emph{Voevodsky motives paper}: \url{https://www.math.ias.edu/Voevodsky/files/files-original/Dropbox/Published_papers/Motives/Collection/s1.pdf}.\item Literature triage checked 2026-08-21, status \texttt{open}; record updated 2026-09-16. The \texttt{corrected\_verified} field reports checking the official report, not independently checked here.\end{itemize}
\end{document}
